\section{Estructura del Sistema}
Esta sección tiene como propósito que el desarrollador comprenda el detalle de implementación de cada clase que compone el sistema, complementando lo ya explicado a nivel conceptual en el Modelo de Dominio (sección 3) y en los Diagramas de Componentes (secciones 9.1 y 9.2). Aquí no se repite el propósito de cada componente, ya justificado en esas secciones, sino que se detallan sus atributos, métodos y asociaciones con el nivel de precisión necesario para guiar la implementación de las clases.

\subsection{Núcleo del dominio}
Transaction es la única clase con identidad propia del dominio y se construye a partir de los objetos de valor Amount, Currency, OrderReference, Payer y GatewayTransactionId, además de referenciar el enum TransactionStatus que indica su estado y, opcionalmente, una instancia de RejectionReason cuando ese estado es DECLINED. Incluye también el atributo opcional authorizationCode, fijado en el constructor cuando la pasarela lo provee. Expone los métodos isApproved(), isFinal() e isPending(), que permiten al código cliente razonar sobre el resultado de un pago sin necesidad de comparar directamente contra los valores del enum. Conserva además el campo rawStatus, que guarda el valor nativo devuelto por la pasarela antes de ser normalizado, con fines de auditoría.

A diferencia de lo que podría sugerir un diagrama con un método de actualización de estado, Transaction es inmutable: ninguno de sus atributos se modifica una vez construida la instancia. Cuando llega un webhook de conciliación con un cambio de estado sobre una transacción que ya existía como PENDING, el ResponseNormalizer no muta la instancia original, sino que construye una nueva instancia de Transaction con el estado actualizado a partir del evento recibido. Esta decisión es consistente con el tratamiento de Transaction como agregado inmutable dentro del dominio y evita efectos secundarios en referencias que otro componente pudiera conservar sobre la instancia previa.

Los objetos de valor que componen a Transaction deben implementarse como inmutables. Amount encapsula un único atributo numérico con hasta dos decimales y expone toMinorUnits() para que cada Adapter lo traduzca a la representación nativa de su pasarela, además de un método equals() para comparación por valor. Currency valida en su constructor que el código recibido tenga exactamente tres letras mayúsculas conforme a ISO 4217, con "COP" como valor por defecto. OrderReference valida únicamente que su valor no esté vacío, sin generar ni garantizar unicidad. Payer agrupa los datos de quien paga y valida en su constructor que el campo email esté presente, siendo el único atributo obligatorio. GatewayTransactionId combina el identificador nativo con el Gateway que lo originó, y se diferencia de OrderReference en que es la pasarela quien lo genera, no el comercio. RejectionReason combina el código nativo del rechazo (rejectionCode) con la categoría normalizada (rejectionCategory) tomada del enum RejectionCategory. ReturnUrlConfig expone el método resolveFor(status), que retorna la URL correspondiente según el resultado del pago, soportando tanto una URL única como URLs diferenciadas por resultado.

WebhookEvent es un objeto de valor inmutable del dominio que representa el evento notificado de forma asíncrona por una pasarela cuando ocurre un cambio de estado en una transacción. Encapsula el identificador de transacción nativo (gatewayTransactionId), el tipo de evento normalizado, la pasarela de origen (gateway), el estado resultante mapeado al dominio (status) y el payload nativo original (rawPayload). Es retornado por las operaciones de parsing de webhooks para permitir una manipulación tipada y uniforme de las notificaciones sin depender de las estructuras heterogéneas de cada proveedor.

WebhookVerifier se implementa como servicio de dominio sin estado propio. Expone dos métodos públicos: verify(payload, headers, secret, gateway): boolean, que valida la integridad criptográfica de la firma del mensaje, y parse(payload, gateway): WebhookEvent, que transforma la notificación entrante de una pasarela específica en una instancia normalizada de WebhookEvent para dar cumplimiento a RF-04. Internamente delega la verificación y extracción de atributos en la lógica correspondiente al Gateway recibido.

SdkError extiende la clase Error nativa de JavaScript y añade los atributos code (del tipo SdkErrorCode), gateway y originalPayload, este último tipado como unknown para forzar una verificación explícita antes de su uso. Los enums TransactionStatus, RejectionCategory, SdkErrorCode y Gateway no requieren métodos propios; su detalle de valores ya está documentado en la sección 3 y no se repite aquí.

\subsection{SDK}
PaymentFacade es la única clase que el desarrollador que consume el SDK instancia directamente. Mantiene una referencia a SdkConfigurator y a GatewayFactory, y expone los tres métodos públicos del SDK: createPayment(), getPaymentStatus() y validateWebhook(). Antes de ejecutar cualquier operación, consulta al Configurator para obtener la pasarela activa y sus credenciales, solicita al Factory la instancia del adaptador correspondiente, y envuelve la llamada resultante con RetryHandler.

SdkConfigurator centraliza el estado de configuración del SDK: la pasarela activa, un mapa de credenciales por proveedor, el entorno de ejecución (production, sandbox o simulation) y los parámetros del mecanismo de reintento. Expone configure(options) para inicializar esos valores y métodos de acceso puntual como getActiveGateway() y getCredentials(gateway). GatewayFactory no mantiene estado; su único método, create(gateway), recibe un valor del enum Gateway y retorna la instancia del Adapter correspondiente, ocultando ese proceso de instanciación del resto del sistema.

Los cuatro Adapters (WompiAdapter, RapydAdapter, MercadoPagoAdapter y KushkiAdapter) implementan la interfaz PaymentGatewayPort, que define los métodos createPayment(), getStatus() y verifySignature(). Cada Adapter mantiene su propio cliente HTTP y sus credenciales, y traduce internamente esos tres métodos hacia las convenciones específicas de su proveedor: autenticación por encabezados HTTP (\texttt{access\_key}, salt, timestamp y firma) en Rapyd, Bearer en Mercado Pago, el header Private-Merchant-Id en Kushki, y el algoritmo de firma correspondiente en cada caso (SHA-256 para Wompi, HMAC-SHA256 con codificación Base64 para Rapyd, HMAC-SHA256 para Mercado Pago y Kushki). Ninguno de los cuatro expone métodos adicionales a los definidos por el puerto, para mantener la sustituibilidad entre ellos.

ResponseNormalizer mantiene internamente la tabla de equivalencias entre los estados nativos de cada pasarela y el enum TransactionStatus. Su método normalize(rawResponse, gateway) toma la respuesta cruda de cualquier Adapter y construye la entidad Transaction correspondiente, siendo invocado por PaymentFacade inmediatamente después de cada llamada exitosa a un Adapter.

RetryHandler envuelve las llamadas que PaymentFacade hace a los Adapters mediante su método execute(operation), que ejecuta la operación recibida y, ante una falla, consulta isTransient(error) para decidir si corresponde reintentar. Mantiene como atributos el número máximo de intentos (tres por defecto) y los intervalos de espera entre ellos (uno, dos y cuatro segundos por defecto). Si se agotan los intentos sin éxito, delega el error a ErrorHandler con el código MAX\_RETRIES\_EXCEEDED.

ErrorHandler mantiene un mapa de errores nativos hacia SdkErrorCode por cada Gateway, y expone el método handle(rawError, gateway), que traduce cualquier error recibido de un Adapter o del RetryHandler hacia una instancia de SdkError, siendo este el único tipo de error que finalmente se propaga al código cliente del SDK.

\subsection{API de Simulación}
HttpRouterMiddleware, implementado como plugin de Fastify, es responsable de dos tareas antes de delegar el procesamiento: identificar la pasarela de destino a partir de la estructura de la URL mediante identifyGateway(url), y extraer los headers de control de la solicitud mediante extractHeaders(req), que retorna el escenario solicitado (\texttt{x-simulate-\allowbreak scenario}) y la latencia a introducir (\texttt{x-simulate-\allowbreak delay}).

ScenarioExecutionEngine concentra la lógica de decisión del simulador en un único método, execute(scenario, gateway). Para los escenarios de negocio (aprobado, rechazado, fondos insuficientes) delega la construcción del payload a GatewayMockFactory; para el escenario de timeout, introduce la latencia recibida antes de responder; para error de red, retorna directamente una respuesta HTTP 500 o 502 sin invocar a la fábrica.

GatewayMockFactory expone el método build(scenario, gateway) y se compone internamente de cuatro fábricas específicas por proveedor: WompiMockFactory, RapydMockFactory, MercadoPagoMockFactory y KushkiMockFactory. Cada una de estas cuatro clases internas conoce el vocabulario de campos y valores propio de su pasarela y expone su propio método build(scenario), sin ser accesibles directamente desde fuera de GatewayMockFactory.

SignatureGenerator expone el método sign(payload, gateway, secret) y es invocado exclusivamente por WebhookTriggerEndpoint, nunca de forma automática. WebhookTriggerEndpoint implementa el endpoint \texttt{POST /v1/sim/\allowbreak webhooks/trigger} mediante su método trigger(targetUrl, eventType, gateway), que solicita la firma correspondiente a SignatureGenerator antes de despachar el webhook sintético hacia la URL indicada por quien lo invoca.

OpenApiDocumentationProvider, implementado mediante el plugin Swagger de Fastify, no expone métodos propios de negocio: genera automáticamente la especificación OpenAPI 3.0 a partir de las definiciones de rutas registradas en HttpRouterMiddleware, sirviéndola en el endpoint /docs.

